\documentclass[10pt,a4paper]{amsart}
\usepackage{amsmath,amssymb,amsthm,graphicx}
\usepackage{verbatim}
\newenvironment{reference}{\comment}{\endcomment}
%\newenvironment{reference}{}{}
\newenvironment{slogan}{\comment}{\endcomment}
\newenvironment{history}{\comment}{\endcomment}

% For commutative diagrams we use Xy-pic
\usepackage[all]{xy}

% We use 2cell for 2-commutative diagrams.
\xyoption{2cell}
% Theorem environments.
%
\theoremstyle{plain}
\newtheorem{theorem}[subsection]{Theorem}
\newtheorem{proposition}[subsection]{Proposition}
\newtheorem{lemma}[subsection]{Lemma}

\theoremstyle{definition}
\newtheorem{definition}[subsection]{Definition}
\newtheorem{example}[subsection]{Example}
\newtheorem{exercise}[subsection]{Exercise}
\newtheorem{situation}[subsection]{Situation}

\theoremstyle{remark}
\newtheorem{remark}[subsection]{Remark}
\newtheorem{remarks}[subsection]{Remarks}

\numberwithin{equation}{subsection}

% Macros
%
\def\lim{\mathop{\mathrm{lim}}\nolimits}
\def\colim{\mathop{\mathrm{colim}}\nolimits}
\def\Spec{\mathop{\mathrm{Spec}}}
\def\Hom{\mathop{\mathrm{Hom}}\nolimits}
\def\Ext{\mathop{\mathrm{Ext}}\nolimits}
\def\SheafHom{\mathop{\mathcal{H}\!\mathit{om}}\nolimits}
\def\SheafExt{\mathop{\mathcal{E}\!\mathit{xt}}\nolimits}
\def\Sch{\mathit{Sch}}
\def\Mor{\mathop{\mathrm{Mor}}\nolimits}
\def\Ob{\mathop{\mathrm{Ob}}\nolimits}
\def\Sh{\mathop{\mathit{Sh}}\nolimits}
\def\NL{\mathop{N\!L}\nolimits}
\def\CH{\mathop{\mathrm{CH}}\nolimits}
\def\proetale{{pro\text{-}\acute{e}tale}}
\def\etale{{\acute{e}tale}}
\def\QCoh{\mathit{QCoh}}
\def\Ker{\mathop{\mathrm{Ker}}}
\def\Im{\mathop{\mathrm{Im}}}
\def\Coker{\mathop{\mathrm{Coker}}}
\def\Coim{\mathop{\mathrm{Coim}}}

% Boxtimes
%
\DeclareMathSymbol{\boxtimes}{\mathbin}{AMSa}{"02}

%
% Macros for moduli stacks/spaces
%
\def\QCohstack{\mathcal{QC}\!\mathit{oh}}
\def\Cohstack{\mathcal{C}\!\mathit{oh}}
\def\Spacesstack{\mathcal{S}\!\mathit{paces}}
\def\Quotfunctor{\mathrm{Quot}}
\def\Hilbfunctor{\mathrm{Hilb}}
\def\Curvesstack{\mathcal{C}\!\mathit{urves}}
\def\Polarizedstack{\mathcal{P}\!\mathit{olarized}}
\def\Complexesstack{\mathcal{C}\!\mathit{omplexes}}
% \Pic is the operator that assigns to X its picard group, usage \Pic(X)
% \Picardstack_{X/B} denotes the Picard stack of X over B
% \Picardfunctor_{X/B} denotes the Picard functor of X over B
\def\Pic{\mathop{\mathrm{Pic}}\nolimits}
\def\Picardstack{\mathcal{P}\!\mathit{ic}}
\def\Picardfunctor{\mathrm{Pic}}
\def\Deformationcategory{\mathcal{D}\!\mathit{ef}}
\usepackage{cleveref}
\crefname{originalreference}{original reference}{original references}
\Crefname{originalreference}{Original reference}{Original references}
\makeatletter
\newlabel{example-map-hom-tensor}{{0BYX}{source}{}{sourcepointer-1}{}}
\newlabel{example-map-hom-tensor@cref}{{[originalreference][0][]0BYX}{[1][0][]source}{}{}{}}
\makeatother
\crefname{originalreference}{original reference}{original references}
\Crefname{originalreference}{Original reference}{Original references}
\makeatletter
\newlabel{lemma-compose-tensor-functors}{{09S4}{source}{}{sourcepointer-2}{}}
\newlabel{lemma-compose-tensor-functors@cref}{{[originalreference][0][]09S4}{[1][0][]source}{}{}{}}
\makeatother
\crefname{originalreference}{original reference}{original references}
\Crefname{originalreference}{Original reference}{Original references}
\makeatletter
\newlabel{lemma-hom-derived}{{09KY}{source}{}{sourcepointer-3}{}}
\newlabel{lemma-hom-derived@cref}{{[originalreference][0][]09KY}{[1][0][]source}{}{}{}}
\makeatother
\crefname{originalreference}{original reference}{original references}
\Crefname{originalreference}{Original reference}{Original references}
\makeatletter
\newlabel{lemma-qis-equivalence}{{09S6}{source}{}{sourcepointer-4}{}}
\newlabel{lemma-qis-equivalence@cref}{{[originalreference][0][]09S6}{[1][0][]source}{}{}{}}
\makeatother
\crefname{originalreference}{original reference}{original references}
\Crefname{originalreference}{Original reference}{Original references}
\makeatletter
\newlabel{lemma-tilting-equivalence}{{09S7}{source}{}{sourcepointer-5}{}}
\newlabel{lemma-tilting-equivalence@cref}{{[originalreference][0][]09S7}{[1][0][]source}{}{}{}}
\makeatother
\crefname{originalreference}{original reference}{original references}
\Crefname{originalreference}{Original reference}{Original references}
\makeatletter
\newlabel{remark-lift-equivalence-to-dga}{{09S9}{source}{}{sourcepointer-6}{}}
\newlabel{remark-lift-equivalence-to-dga@cref}{{[originalreference][0][]09S9}{[1][0][]source}{}{}{}}
\makeatother
\title{A linear equivalence of derived categories with concentrated source and K-flat target comes from a bimodule}
\author{The Stacks Project Authors}
\date{stacks-project:dga.tex:local-snapshot-20261012; GNU FDL 1.2 or later}
\begin{document}
\maketitle
\noindent\textit{Editorial note.} All delivered mathematical text is the author's own, copied from the pinned Stacks Project chapter file sources/stacks-dga/pinned-original-source.tex; no statement or proof step was rewritten, reordered or removed.\par\medskip
\noindent\textit{Editorial note.} The counted proof is the author's complete argument: the choice N = F(A) as a property-(P) module, the endomorphism dga E with H\textasciicircum{}0(E) = A, the degree-wise truncated subalgebra E' with its quasi-isomorphism zigzag, the equivalence chain with the quasi-inverse identification M maps to M tensor\_A\textasciicircum{}L A, and the closing composition through lemma-compose-tensor-functors are all retained verbatim.\par\medskip
\noindent\textit{Editorial note.} The label proposition-rickard names the counted statement, and the proof environment immediately following it in the source is delivered verbatim as its proof.\par\medskip
\noindent\textit{Editorial note.} The counted statement is carried by the author's own proposition environment in the chapter source.\par\medskip
\noindent\textit{Allowed standard prerequisites (editorial).}\begin{itemize}
\item The endomorphism differential graded algebra of the image of the unit. Same chapter, Remark remark-lift-equivalence-to-dga, Lemma lemma-tilting-equivalence and Lemma lemma-hom-derived. Sets N = F(A) as a property-(P) module, identifies (E, d) = Hom(N, N) with H\textasciicircum{}0(E) = A and vanishing higher cohomology, and gives through the tilting lemma that N as an (E, B)-bimodule induces an equivalence D(E, d) to D(B, d).
\item Derived categories of quasi-isomorphic algebras and restriction of scalars. Same chapter, Lemma lemma-qis-equivalence and Example example-map-hom-tensor. Turns the quasi-isomorphism zigzag of algebras into equivalences of derived categories and identifies the quasi-inverse of restriction with derived tensor.
\end{itemize}\medskip
\ifdefined\hypertarget\hypertarget{sourcepointer-1}{}\fi
\noindent\textit{Source reference.} Original reference 0BYX: sources/stacks-dga/pinned-original-source.tex, source lines 5657--5657. Prerequisite pointer invoked at proof line(s) 6950; it names a standard fact of the same chapter that the counted argument uses, and every step of the counted proof built on it is retained verbatim in the delivered body.\par\medskip
\ifdefined\hypertarget\hypertarget{sourcepointer-2}{}\fi
\noindent\textit{Source reference.} Original reference 09S4: sources/stacks-dga/pinned-original-source.tex, source lines 6067--6067. Prerequisite pointer invoked at proof line(s) 6953; it names a standard fact of the same chapter that the counted argument uses, and every step of the counted proof built on it is retained verbatim in the delivered body.\par\medskip
\ifdefined\hypertarget\hypertarget{sourcepointer-3}{}\fi
\noindent\textit{Source reference.} Original reference 09KY: sources/stacks-dga/pinned-original-source.tex, source lines 2850--2850. Prerequisite pointer invoked at proof line(s) 6923; it names a standard fact of the same chapter that the counted argument uses, and every step of the counted proof built on it is retained verbatim in the delivered body.\par\medskip
\ifdefined\hypertarget\hypertarget{sourcepointer-4}{}\fi
\noindent\textit{Source reference.} Original reference 09S6: sources/stacks-dga/pinned-original-source.tex, source lines 6725--6725. Prerequisite pointer invoked at proof line(s) 6946; it names a standard fact of the same chapter that the counted argument uses, and every step of the counted proof built on it is retained verbatim in the delivered body.\par\medskip
\ifdefined\hypertarget\hypertarget{sourcepointer-5}{}\fi
\noindent\textit{Source reference.} Original reference 09S7: sources/stacks-dga/pinned-original-source.tex, source lines 6753--6753. Prerequisite pointer invoked at proof line(s) 6912; it names a standard fact of the same chapter that the counted argument uses, and every step of the counted proof built on it is retained verbatim in the delivered body.\par\medskip
\ifdefined\hypertarget\hypertarget{sourcepointer-6}{}\fi
\noindent\textit{Source reference.} Original reference 09S9: sources/stacks-dga/pinned-original-source.tex, source lines 6865--6865. Prerequisite pointer invoked at proof line(s) 6916; it names a standard fact of the same chapter that the counted argument uses, and every step of the counted proof built on it is retained verbatim in the delivered body.\par\medskip
\begin{proposition}
\label{proposition-rickard}
Let $R$ be a ring. Let $(A, \text{d})$ and $(B, \text{d})$ be
differential graded $R$-algebras. Let $F : D(A, \text{d}) \to D(B, \text{d})$
be an $R$-linear equivalence of triangulated categories. Assume that
\begin{enumerate}
\item $A = H^0(A)$, and
\item $B$ is K-flat as a complex of $R$-modules.
\end{enumerate}
Then there exists an $(A, B)$-bimodule $N$ as in
Lemma \ref{lemma-tilting-equivalence}.
\end{proposition}

\begin{proof}
As in Remark \ref{remark-lift-equivalence-to-dga} above, we set $N = F(A)$
in $D(B, \text{d})$. We may assume that $N$ is a differential graded
$B$-module with property (P). Set
$$
(E, \text{d}) = \Hom_{\text{Mod}^{dg}_{(B, \text{d})}}(N, N)
$$
Then $H^0(E) = A$ and $H^k(E) = 0$ for $k \not = 0$ by
Lemma \ref{lemma-hom-derived}.
Moreover, by the discussion in Remark \ref{remark-lift-equivalence-to-dga}
and by Lemma \ref{lemma-tilting-equivalence}
we see that $N$ as a $(E, B)$-bimodule induces an
equivalence $- \otimes_E^\mathbf{L} N : D(E, \text{d}) \to D(B, \text{d})$.
Let $E' \subset E$ be the differential graded $R$-subalgebra
with
$$
(E')^i = \left\{
\begin{matrix}
E^i & \text{if }i < 0 \\
\Ker(E^0 \to E^1) & \text{if }i = 0 \\
0 & \text{if }i > 0
\end{matrix}
\right.
$$
Then there are quasi-isomorphisms of differential graded
algebras $(A, \text{d}) \leftarrow (E', \text{d}) \rightarrow (E, \text{d})$.
Thus we obtain equivalences
$$
D(A, \text{d}) \leftarrow D(E', \text{d}) \rightarrow D(E, \text{d})
\rightarrow D(B, \text{d})
$$
by Lemma \ref{lemma-qis-equivalence}.
Note that the quasi-inverse $D(A, \text{d}) \to D(E', \text{d})$
of the left vertical arrow is given
by $M \mapsto M \otimes_A^\mathbf{L} A$ where $A$ is viewed as a
$(A, E')$-bimodule, see Example \ref{example-map-hom-tensor}.
On the other hand the functor $D(E', \text{d}) \to D(B, \text{d})$ is given by
$M \mapsto M \otimes_{E'}^\mathbf{L} N$ where $N$ is as above.
We conclude by Lemma \ref{lemma-compose-tensor-functors}.
\end{proof}

\end{document}
